% ===== Block B00a: 原文 p.v (Preface) =====
% [原文 p.v]
% [Tutor Note: 前言给出全书的几何路线图：Ch1 曲面/曲线、Ch2 仿射李代数、
%  Ch3 构造共形块、Ch4 层的构造、Ch5 射影平坦联络、Ch6 球面情形。
%  精读时建议先读本页，把每章的"输入—输出"记下来。]

\begin{sloppypar}本讲义的主要目的是从复代数几何的视角，对带规范对称性的共形场论——即所谓的
Wess--Zumino--Witten--Novikov 模型——给出一套系统的处理方式。在第 1 章与第 2 章中
分别介绍紧黎曼曲面理论的基本事实与仿射李代数的表示论之后，我们将在第 3 章对带坐标
的稳定带点曲线构造共形块。在第 4 章中，我们将构造与带坐标的稳定带点曲线族相伴的
共形块层。第 5 章将证明共形块层带有射影平坦联络，这是共形场论中最重要的事实之一。
第 6 章致力于研究 $\PP^1$ 上共形场论的细致结构。
\end{sloppypar}

% [Tutor Note: "with coordinates" 是全书的要害：共形块依赖于每个 marked point 处
%  的形式坐标（或局部参数）。第 3 章的定义与第 4 章的层化都以此为出发点。]

最近，J.\,E.\ Andersen 与我从共形场论构造了模函子（modular functor）。这给出了
代数几何与拓扑量子场论之间的一个有趣联系。本讲义包含了构造该模函子所用的、关于带
规范对称性共形场论的全部必要技巧与结果。

% [Tutor Note: modular functor 一句解释了本书的动机：Ch1--Ch6 的技术最终服务于
%  从共形场论构造 modular functor（即与曲面模空间相容的、满足拼接公理的向量空间系）。
%  若只关心代数几何侧，可把 Ch2、Ch6 暂时跳读。]

本讲义基于在 Fields Institute、Queen's University、Århus University、Kobe
University 与 Kyoto University 所做的一系列讲座与报告。我感谢这些热情的听众，
他们帮助我改进了本讲义中的某些部分。

我要感谢 Noriko Yui 邀请我为 Fields Institute Monographs 撰写本讲义，并帮助我
完成手稿。最后但同样重要的是，我要向 Mathematical Reviews 的 Arthur Greenspoon
表达衷心的谢意，他润色了我原稿的英文。

\vspace{1.2em}
\begin{flushright}
2008 年 3 月\\
Kenji Ueno
\end{flushright}
